% ═══════════════════════════════════════════════════════════════════════════
% LEHRERHINWEISE: Me gusta - Vorlieben ausdrücken
% Fach: Spanisch (Oberstufe Spätbeginner)
% Kurz-Lessonplan + Lösungen
% ═══════════════════════════════════════════════════════════════════════════

\documentclass[10pt, a4paper]{article}

% ═══════════════════════════════════════════════════════════════════════════
% PAKETE
% ═══════════════════════════════════════════════════════════════════════════

\usepackage[utf8]{inputenc}
\usepackage[T1]{fontenc}
\usepackage[ngerman]{babel}
\usepackage{helvet}
\renewcommand{\familydefault}{\sfdefault}

\usepackage[margin=1.8cm, top=2cm, bottom=1.8cm]{geometry}
\usepackage{enumitem}
\usepackage{tabularx}
\usepackage{booktabs}
\usepackage{array}
\usepackage{multirow}
\usepackage{multicol}

\usepackage{xcolor}
\usepackage{tcolorbox}
\usepackage{fancyhdr}
\usepackage{lastpage}
\usepackage{titlesec}

% Kompaktere Absätze
\setlength{\parskip}{4pt}
\setlength{\parindent}{0pt}

% ═══════════════════════════════════════════════════════════════════════════
% FARBEN
% ═══════════════════════════════════════════════════════════════════════════

\definecolor{spanisch}{HTML}{c41e3a}
\definecolor{grau}{HTML}{666666}
\definecolor{hellgrau}{HTML}{f5f5f5}
\definecolor{success}{HTML}{2a7a4b}
\definecolor{warning}{HTML}{b35c00}
\definecolor{fehler}{HTML}{c62828}

% ═══════════════════════════════════════════════════════════════════════════
% ÜBERSCHRIFTEN
% ═══════════════════════════════════════════════════════════════════════════

\titleformat{\section}
    {\normalfont\large\bfseries\color{spanisch}}
    {\thesection}{0.8em}{}

\titleformat{\subsection}
    {\normalfont\normalsize\bfseries}
    {\thesubsection}{0.8em}{}

\titlespacing*{\section}{0pt}{12pt}{6pt}
\titlespacing*{\subsection}{0pt}{8pt}{4pt}

% ═══════════════════════════════════════════════════════════════════════════
% KOPF- UND FUSSZEILE
% ═══════════════════════════════════════════════════════════════════════════

\pagestyle{fancy}
\fancyhf{}
\renewcommand{\headrulewidth}{0.4pt}
\renewcommand{\footrulewidth}{0pt}

\lhead{\small\textbf{Lehrerhinweise} -- Spanisch}
\rhead{\small Me gusta -- Stunde 1/3}
\cfoot{\small Seite \thepage\ von \pageref{LastPage}}

% ═══════════════════════════════════════════════════════════════════════════
% BEFEHLE
% ═══════════════════════════════════════════════════════════════════════════

\newcommand{\es}[1]{\textcolor{spanisch}{\textit{#1}}}
\newcommand{\loesung}[1]{\textcolor{success}{\textbf{#1}}}

% Tipp-Box
\newtcolorbox{tippbox}{
    colback=hellgrau,
    colframe=spanisch,
    boxrule=0.8pt,
    arc=0pt,
    left=6pt, right=6pt, top=4pt, bottom=4pt,
    title={\small\bfseries Tipp}
}

% Warnung-Box
\newtcolorbox{warnbox}{
    colback=white,
    colframe=warning,
    boxrule=1pt,
    arc=0pt,
    left=6pt, right=6pt, top=4pt, bottom=4pt,
    title={\small\bfseries Achtung}
}

% Fehler-Box
\newtcolorbox{fehlerbox}{
    colback=white,
    colframe=fehler,
    boxrule=0.8pt,
    arc=0pt,
    left=6pt, right=6pt, top=4pt, bottom=4pt
}

% ═══════════════════════════════════════════════════════════════════════════
% DOKUMENT
% ═══════════════════════════════════════════════════════════════════════════

\begin{document}

% ═══════════════════════════════════════════════════════════════════════════
% TITEL
% ═══════════════════════════════════════════════════════════════════════════

\begin{center}
    {\Large\bfseries\color{spanisch} Lehrerhinweise: Me gusta}\\[4pt]
    {\normalsize Stunde 1 von 3 -- Einführung \es{gustar} + Hobbys}\\[2pt]
    {\small\color{grau} Oberstufe Spätbeginner | 45 Minuten | GER A2}
\end{center}

\vspace{-4pt}
\hrule
\vspace{8pt}

% ═══════════════════════════════════════════════════════════════════════════
% KURZÜBERSICHT
% ═══════════════════════════════════════════════════════════════════════════

\begin{multicols}{2}

\section*{Kurzübersicht}

\textbf{Lernziele:}
\begin{itemize}[nosep, leftmargin=1.5em]
    \item Struktur von \es{gustar} verstehen (= gefallen)
    \item Pronomen me/te/le/nos/os/les anwenden
    \item gusta (Sg./Inf.) vs. gustan (Pl.) unterscheiden
    \item Vorlieben ausdrücken und erfragen
\end{itemize}

\textbf{Material:}
\begin{itemize}[nosep, leftmargin=1.5em]
    \item[$\Box$] Beamer + Einstiegsfolie
    \item[$\Box$] Tablets/PCs (oder nur analog)
    \item[$\Box$] Lernpfad: \texttt{lernpfad.html}
    \item[$\Box$] AB kopiert (1x pro SuS)
\end{itemize}

\columnbreak

\textbf{Vorbereitung (5 min):}
\begin{enumerate}[nosep, leftmargin=1.5em]
    \item Einstiegsfolie am Beamer bereit
    \item Lernpfad-URL an Tafel/Board
    \item AB-Stapel bereitlegen
    \item Ggf. Tablets verteilen lassen
\end{enumerate}

\textbf{Kernbotschaft der Stunde:}

\begin{tippbox}
\es{Gustar} funktioniert wie \textbf{„gefallen"}, nicht wie „mögen"!\\
Das Subjekt ist das, WAS gefällt -- nicht die Person.
\end{tippbox}

\end{multicols}

\vspace{-8pt}
\hrule
\vspace{6pt}

% ═══════════════════════════════════════════════════════════════════════════
% STUNDENVERLAUF
% ═══════════════════════════════════════════════════════════════════════════

\section{Stundenverlauf (5-Minuten-Raster)}

\small
\begin{tabularx}{\textwidth}{|c|c|X|X|c|}
    \hline
    \textbf{Min.} & \textbf{Phase} & \textbf{Inhalt} & \textbf{L-Aktivität} & \textbf{SF} \\
    \hline
    0-5 & Einstieg & „Überraschungsfakt": Du kannst nicht sagen „Ich mag dich" & Folie zeigen, Frage stellen: „Warum?" & PL \\
    \hline
    5-10 & Erarbeitung & Lernpfad Schritt 1-2: \es{gustar} = gefallen & AA geben, herumgehen, Fragen beantworten & EA \\
    \hline
    10-15 & Übung I & Lernpfad Schritt 3: Verständnis-MC & Fortschritt beobachten & EA \\
    \hline
    15-20 & Erarbeitung & Lernpfad Schritt 4: Pronomen + AB Kasten 1-2 & AB austeilen, Übertragen anleiten & EA \\
    \hline
    20-25 & Übung II & Lernpfad Schritt 5/5b: Pronomen + gusta/gustan Drill & Beobachten, bei Bedarf Hilfestellung & EA \\
    \hline
    25-30 & Vokabeln & Lernpfad Schritt 6: Hobbys & Kurze Besprechung im Plenum möglich & EA/PL \\
    \hline
    30-35 & Anwendung & AB Ü1-3 in Partnerarbeit & Herumgehen, Peer-Korrektur fördern & PA \\
    \hline
    35-40 & Test & Lernpfad Schritt 8: Abschlusstest (6 MC) & Ruhe, keine Hilfe & EA \\
    \hline
    40-45 & Sicherung & AB Ü4 (Fehler erklären) + Freies Schreiben & Einsammeln, Ausblick Stunde 2 & EA \\
    \hline
\end{tabularx}
\normalsize

\textbf{Legende:} PL = Plenum, EA = Einzelarbeit, PA = Partnerarbeit, L = Lehrkraft, AA = Arbeitsauftrag

\vspace{4pt}

\begin{warnbox}
\textbf{Zeitpuffer:} Schritt 5b (Drill) kann bei Zeitnot gekürzt werden. Bei schneller Gruppe: Partner-Interview ergänzen (\es{¿Qué te gusta?}).
\end{warnbox}

% ═══════════════════════════════════════════════════════════════════════════
% DIFFERENZIERUNG
% ═══════════════════════════════════════════════════════════════════════════

\section{Differenzierung}

\begin{tabularx}{\textwidth}{|c|X|X|}
    \hline
    \textbf{Niveau} & \textbf{Aufgaben} & \textbf{Unterstützung / Erweiterung} \\
    \hline
    \textbf{Basis} (*) & AB Ü1 (Pronomen einsetzen), AB Ü2 (gusta/gustan), Lernpfad Schritt 3, 5 &
    Pronomen-Tabelle als Hilfe bereitstellen; mehr Zeit; Partnerarbeit mit stärkerem SuS \\
    \hline
    \textbf{Standard} (**) & AB Ü3 (Übersetzen), Lernpfad Schritt 7 &
    Selbstständig arbeiten; Peer-Korrektur \\
    \hline
    \textbf{Erweitert} (***) & AB Ü4 (Fehler begründen), Freies Schreiben, ggf. Partner-Interview &
    Zusatzaufgabe: weitere Verben (\es{encantar, interesar}); Expertenrolle bei PA \\
    \hline
\end{tabularx}

\vspace{6pt}

\textbf{Individuelle Förderung:}
\begin{itemize}[nosep, leftmargin=1.5em]
    \item \textbf{LRS:} AB in Schriftgröße 12pt, mehr Zeit, ggf. mündlich prüfen
    \item \textbf{Vorkenntnisse:} Als Experte einsetzen, Zusatzverben, eigene Beispielsätze
    \item \textbf{Zurückhaltende SuS:} Gezielt in PA mit vertrautem Partner; erst schriftlich, dann mündlich
\end{itemize}

% ═══════════════════════════════════════════════════════════════════════════
% HÄUFIGE FEHLER
% ═══════════════════════════════════════════════════════════════════════════

\section{Häufige Schülerfehler und Reaktionen}

\begin{tabularx}{\textwidth}{|X|X|X|}
    \hline
    \textbf{Fehler} & \textbf{Ursache} & \textbf{Reaktion / Korrektur} \\
    \hline
    \es{*Yo gusto el chocolate.} &
    Transfer aus dem Englischen („I like...") &
    Zurück zur Brücke: „Wie heißt das auf Deutsch mit ‚gefallen'?" → \es{Me gusta el chocolate.} \\
    \hline
    \es{*Me gusta los libros.} &
    Plural nicht erkannt → falsches Verb &
    Frage: „Was ist das Subjekt? Singular oder Plural?" → \es{gustan} \\
    \hline
    \es{*Me gustan bailar.} &
    Infinitiv als Plural interpretiert &
    Erklären: „Infinitiv = eine Handlung = Singular" → \es{gusta} \\
    \hline
    \es{*Mi gusta...} statt \es{Me gusta...} &
    Verwechslung mit Possessivpronomen \es{mi} &
    Unterschied klären: \es{mi} (mein) $\neq$ \es{me} (mir) \\
    \hline
    \es{*A mí gusta...} (ohne \es{me}) &
    Glaubt, \es{a mí} ersetzt \es{me} &
    Regel: \es{me/te/le...} ist IMMER nötig, \es{a mí/a ti...} ist optional \\
    \hline
    Artikel vergessen: \es{*Me gusta música.} &
    Im Deutschen ohne Artikel („Mir gefällt Musik") &
    Im Spanischen braucht man den Artikel: \es{la música} \\
    \hline
\end{tabularx}

\vspace{6pt}

\begin{tippbox}
\textbf{Korrektives Feedback:} Nicht direkt „falsch" sagen, sondern nachfragen: „Was ist das Subjekt?" / „Singular oder Plural?" / „Wie sagt man das mit ‚gefallen'?" -- So aktivieren SuS ihr Wissen selbst.
\end{tippbox}

% ═══════════════════════════════════════════════════════════════════════════
% LÖSUNGEN
% ═══════════════════════════════════════════════════════════════════════════

\section{Lösungen zum Arbeitsblatt}

\subsection*{Übung 1: Pronomen einsetzen}

\begin{enumerate}[nosep, leftmargin=2em]
    \item \loesung{Me} gusta el chocolate. (mir)
    \item \loesung{Te} gustan los videojuegos. (dir)
    \item A María \loesung{le} gusta leer. (ihr)
    \item \loesung{Nos} gusta viajar. (uns)
    \item A Pedro y Ana \loesung{les} gustan los gatos. (ihnen)
\end{enumerate}

\subsection*{Übung 2: gusta oder gustan?}

\begin{enumerate}[nosep, leftmargin=2em]
    \item Me \loesung{gusta} la música. (Singular)
    \item Te \loesung{gustan} los deportes. (Plural)
    \item Le \loesung{gusta} bailar. (Infinitiv → Singular)
    \item Nos \loesung{gustan} las películas. (Plural)
    \item Les \loesung{gusta} el fútbol. (Singular)
\end{enumerate}

\subsection*{Übung 3: Übersetze ins Spanische}

\begin{enumerate}[nosep, leftmargin=2em]
    \item Mir gefällt Musik. → \loesung{Me gusta la música.}
    \item Dir gefallen die Filme. → \loesung{Te gustan las películas.}
    \item Uns gefällt das Reisen. → \loesung{Nos gusta viajar.}
\end{enumerate}

\subsection*{Übung 4: Fehler erklären}

\begin{enumerate}[nosep, leftmargin=2em]
    \item \textcolor{fehler}{*Me gustan bailar.}

    \loesung{Begründung:} „Bailar" ist ein Infinitiv. Infinitive gelten als Singular, daher muss es \es{gusta} heißen, nicht \es{gustan}.

    \item \textcolor{fehler}{*Yo gusto el chocolate.}

    \loesung{Begründung:} \es{Gustar} funktioniert wie „gefallen", nicht wie „mögen". Das Subjekt ist „el chocolate" (die Schokolade gefällt mir), nicht „yo". Richtig: \es{Me gusta el chocolate.}
\end{enumerate}

\subsection*{Freies Schreiben: Beispiellösungen}

\begin{itemize}[nosep, leftmargin=2em]
    \item \es{Me gusta escuchar música.} (Mir gefällt es, Musik zu hören.)
    \item \es{No me gusta cocinar.} (Mir gefällt Kochen nicht.)
    \item \es{Me gustan los deportes.} (Mir gefallen die Sportarten.)
    \item \es{Me gusta salir con amigos.} (Mir gefällt es, mit Freunden auszugehen.)
\end{itemize}

% ═══════════════════════════════════════════════════════════════════════════
% LÖSUNGEN ABSCHLUSSTEST
% ═══════════════════════════════════════════════════════════════════════════

\section{Lösungen Abschlusstest (Lernpfad)}

\begin{tabularx}{\textwidth}{|c|X|c|}
    \hline
    \textbf{Nr.} & \textbf{Frage} & \textbf{Lösung} \\
    \hline
    1 & Wie funktioniert „gustar"? & \loesung{B) Wie „gefallen"} \\
    \hline
    2 & \_\_\_ gusta el chocolate. (mir) & \loesung{C) Me} \\
    \hline
    3 & Le \_\_\_ los videojuegos. & \loesung{B) gustan} \\
    \hline
    4 & ¿\_\_\_ gusta bailar? (dir) & \loesung{C) Te} \\
    \hline
    5 & Übersetze: Uns gefällt der Sport. & \loesung{C) Nos gusta el deporte.} \\
    \hline
    6 & Was ist das Subjekt in „Me gustan las manzanas"? & \loesung{C) Las manzanas} \\
    \hline
\end{tabularx}

\textbf{Bewertung:} 6/6 = sehr gut | 5/6 = gut | 4/6 = bestanden | <4 = Wiederholung nötig

% ═══════════════════════════════════════════════════════════════════════════
% AUSBLICK + NOTIZEN
% ═══════════════════════════════════════════════════════════════════════════

\section{Ausblick Stunde 2}

\textbf{Thema:} Vertiefung -- weitere Verben mit gleicher Struktur

\textbf{Inhalte:}
\begin{itemize}[nosep, leftmargin=1.5em]
    \item \es{encantar} (begeistern): \es{Me encanta viajar.}
    \item \es{interesar} (interessieren): \es{¿Te interesa el arte?}
    \item \es{molestar} (stören): \es{Me molesta el ruido.}
    \item Dialogisches Sprechen: Partner-Interview über Vorlieben
\end{itemize}

\textbf{Hausaufgabe (optional):}
\begin{itemize}[nosep, leftmargin=1.5em]
    \item 5 Sätze schreiben: Was gefällt dir? Was gefällt dir nicht?
    \item Vokabeln Hobbys lernen
\end{itemize}

\vspace{8pt}
\hrule
\vspace{6pt}

\section*{Notizen zur Stunde}

\vspace{4cm}

\hrule

\vfill

{\small\color{grau}
Lehrerhinweise erstellt für Greenhouse Schools Rostock |
Begleitend zu: LE-01-gustar-langentwurf.pdf, lernpfad.html, arbeitsblatt-gustos.pdf
}

\end{document}
